\section{Proposed MBA-RL for Charging Control}

Each cell is treated as an agent that observes its local voltage, temperature, and SOC and regulates its own charging current. All agents share the objectives of fast charging and SOC balancing across the battery pack. Since each agent makes decisions from local observations during execution while the charging currents must be coordinated toward a common objective, QMIX \cite{26} is adopted for centralized training and decentralized execution. The resulting MBA-RL framework combines the QMIX controller with the cell models introduced in Section II.

\begin{figure}[t]
    \centering
    \includegraphics[width=1\linewidth]{marl_control_framework.png}
    \caption{MARL formulation of the battery charging control problem. Each agent regulates the charging current of one cell while sharing the objectives of fast charging and SOC balancing.}
    \label{fig:MARL_Schematic}
\end{figure}

\subsection{MBA-RL Framework}

The MBA-RL architecture is shown in Fig.~\ref{fig:schematic_method}. All agents share one recurrent Q-network. For agent $i$, the network takes the current observation $o_t^i$, the previous executed action $\tilde{a}_{t-1}^i$, and the agent identifier as inputs and estimates the local action value $Q_i$. \newrev{During centralized training, a mixing network combines the local action values into the joint action value $Q_{\mathrm{tot}}$, with its weights generated by hypernetworks from the global state. The mixing network is not required for action selection during decentralized execution.} Target agent and mixing networks are used to construct the temporal-difference targets during training.

\begin{figure*}[t]
    \centering
    \includegraphics[width=0.9\textwidth,height=\textheight,keepaspectratio]{mba_rl_architecture.png}
    \caption{Architecture of the MBA-RL framework. The shared agent network consists of two multilayer perceptrons (MLPs) and a gated recurrent unit (GRU).}
    \label{fig:schematic_method}
\end{figure*}

The charging control problem is formulated as a Dec-POMDP \cite{27}:

\[
G=
\left\langle
\mathcal{N},
\mathcal{S},
\{\mathcal{A}_i\}_{i=1}^{n},
\mathcal{P},
\mathcal{R},
\{\mathcal{O}_i\}_{i=1}^{n},
\mathcal{O},
\gamma
\right\rangle .
\]

Here, $\mathcal{N}=\{1,\ldots,n\}$ is the set of agents, $\mathcal{S}$ is the global state space, $\mathcal{A}_i$ and $\mathcal{O}_i$ are the action and observation spaces of agent $i$, $\mathcal{P}$ is the transition kernel, $\mathcal{R}$ is the shared reward function, $\mathcal{O}$ is the observation function, and $\gamma$ is the discount factor.

\subsubsection{State and Observation Spaces}

The global state at time step $t$ is formed by the observations of all cells:

\begin{equation}
\label{eq:global_state}
\mathbf{s}_t
=
\left[
o_t^1,o_t^2,\ldots,o_t^n
\right]^{\mathrm{T}}.
\end{equation}

For agent $i$, the local observation is
$o_t^i=[V_i(t),T_i(t),SOC_i(t)]$, where $V_i(t)$, $T_i(t)$, and $SOC_i(t)$ denote the voltage, temperature, and SOC of cell $i$, respectively. Each agent therefore makes its charging decision from local cell information, while the complete state is available during centralized training.

\subsubsection{Action Space}

Each agent selects a requested charging current from
$\mathcal{A}_{\mathrm{cand}}=\{0,0.5,\ldots,7.5\}\,\mathrm{A}$.
Before action selection, an availability mask removes actions that conflict with the current cell state. When $SOC_i\geq0.95$, only zero current is available. When $V_i\geq V_{\max}$ or $T_i\geq T_{\max}$, currents above 1 A are excluded.

The requested joint action and the corresponding available action space are

\begin{equation}
\label{eq:joint_action}
\mathbf{a}_t
=
\left[
a_t^1,a_t^2,\ldots,a_t^n
\right]^{\mathrm{T}},
\qquad
\mathcal{A}(\mathbf{s}_t)
=
\prod_{i=1}^{n}\mathcal{A}_i(\mathbf{s}_t),
\end{equation}

where $\mathcal{A}_i(\mathbf{s}_t)\subseteq\mathcal{A}_{\mathrm{cand}}$. Let $\mathbf{m}_t$ denote the corresponding availability masks. Since the mask only uses the current state, it does not determine whether a requested current will remain feasible over the next decision interval. The requested current is therefore checked by the model-based safety mechanism before execution.

\subsubsection{Reward Function}

The agents share a reward consisting of a time penalty, an SOC balancing penalty, and penalties for voltage and temperature violations:

\begin{equation}
\label{eq:reward}
r_t
=
r_{\mathrm{time}}
+
r_{\mathrm{bal}}
+
r_{\mathrm{safety}}.
\end{equation}

Before the terminal condition is reached, the time penalty is $r_{\mathrm{time}}=-\lambda_{\mathrm{time}}$. The balancing term is

\begin{equation}
\label{eq:balance_reward}
r_{\mathrm{bal}}
=
\begin{cases}
-\lambda_{\mathrm{bal}}(\sigma(t)-\beta),
& \sigma(t)\geq\beta,\\
0,
& \text{otherwise},
\end{cases}
\end{equation}

\newrev{where $\sigma(t)$ is defined in \eqref{eq:soc_standard_deviation} and $\beta$ is the balancing threshold.}

The safety term is
$r_{\mathrm{safety}}=r_{\mathrm{volt}}+r_{\mathrm{temp}}$,
where
$r_{\mathrm{volt}}=\sum_{i=1}^{n}\delta_{v,i}$ and
$r_{\mathrm{temp}}=\sum_{i=1}^{n}\delta_{t,i}$.
The cell-level penalties are defined as

\begin{equation}
\label{eq:safety_reward}
\begin{aligned}
\delta_{v,i}
&=
\begin{cases}
-\lambda_v\left(V_i(t)-V_{\max}\right),
& V_i(t)\geq V_{\max},\\
0,
& \text{otherwise},
\end{cases}
\\[4pt]
\delta_{t,i}
&=
\begin{cases}
-\lambda_T\left(T_i(t)-T_{\max}\right),
& T_i(t)\geq T_{\max},\\
0,
& \text{otherwise}.
\end{cases}
\end{aligned}
\end{equation}

\subsubsection{Transition and Termination}

\newrev{The policy selects a requested current, and the model-based safety mechanism determines the executed current.} Their joint forms are denoted by $\mathbf{a}_t$ and $\tilde{\mathbf{a}}_t$, respectively.

The transition to the next state is determined by the executed currents and the battery dynamics. \newrev{The terminal flag $d_t$ records whether the episode has ended. In the evaluation, reaching the target SOC and balancing conditions ends the charging phase; task success additionally requires the subsequent holding interval to satisfy the target SOC, balancing, and safety conditions. An episode can also end without completing the task, including when the model-based safety mechanism finds no feasible current.} Both requested and executed actions are retained during training so that interventions by the safety mechanism are represented explicitly.

\subsection{Cooperative Value Learning}

MBA-RL uses centralized training to coordinate the charging decisions of different cells. Agent $i$ estimates the local action value $Q_i(\tau_t^i,a_t^i|\theta)$ from its action-observation history $\tau_t^i$. During data collection, each agent selects a requested action using an $\epsilon$-greedy policy:

\begin{equation}
\label{epsilon-greedy}
a_t^i =
\begin{cases}
\text{random available action},
& \text{with probability }\epsilon,\\
\displaystyle
\arg\max_{a\in\mathcal{A}_i(\mathbf{s}_t)}
Q_i(\tau_t^i,a|\theta),
& \text{with probability }1-\epsilon.
\end{cases}
\end{equation}

All agents share the same network parameters $\theta$. A one-hot agent identifier is included in the network input so that the shared network can distinguish among cells.

\newrev{During centralized training, the mixing network combines the local action values into the joint action value}

\begin{equation}
\label{Q-value}
\begin{aligned}
&Q_{\mathrm{tot}}
(\boldsymbol{\tau}_t,\mathbf{s}_t,\mathbf{a}_t|\theta,\phi)\\
&\qquad=f_{\phi}(\mathbf{Q}_t,\mathbf{s}_t),
\end{aligned}
\end{equation}

\newrev{where $\boldsymbol{\tau}_t$ denotes the joint action-observation history and $\mathbf{Q}_t$ is the vector of local values $Q_i(\tau_t^i,a_t^i|\theta)$. The parameters $\theta$ belong to the shared agent network, while $\phi$ includes the mixing network and its state-dependent hypernetworks.}

QMIX imposes the monotonicity constraint

\begin{equation}
\label{QMIX_constraints}
\frac{\partial Q_{\mathrm{tot}}}{\partial Q_i}
\geq 0,
\qquad
i=1,2,\ldots,n.
\end{equation}

This constraint makes decentralized greedy action selection consistent with maximizing the joint action value, allowing each agent to select its action independently during execution.

\newrev{Every $N_t$ training episodes, the online parameters are copied to the target networks: $\theta'\leftarrow\theta$ and $\phi'\leftarrow\phi$.} Considering the terminal flag and the available actions at the next state, the target value is

% Figure artwork still labels target updates as "Soft update"; the text and
% algorithm specify periodic parameter copying. The figure is unchanged.

\begin{equation}
\label{target Q}
\begin{aligned}
y_t^{\mathrm{tot}}
&=r_t+\gamma(1-d_t)\\
&\quad\times\max_{\mathbf{a}\in\mathcal{A}(\mathbf{s}_{t+1})}
Q_{\mathrm{tot}}
(\boldsymbol{\tau}_{t+1},\mathbf{s}_{t+1},\mathbf{a}|\theta',\phi').
\end{aligned}
\end{equation}

The shared reward couples the local charging decisions through $Q_{\mathrm{tot}}$. The time penalty encourages faster charging, while the balancing term penalizes SOC differences among cells. The agents are therefore trained to coordinate their charging currents toward the common objectives of fast charging and SOC balancing.

\subsection{Model-Based Safety Mechanism}

At each decision step, MBA-RL first generates a requested current for each cell. \newrev{Starting from this value, the model-based safety mechanism evaluates the requested current and successively lower candidates using the nominal SPMe over the next decision interval, followed by an interval with zero current.}

A correction based on the mismatch between the model prediction and the observed cell response is applied to the predicted peaks of SOC, voltage, and temperature. \newrev{The voltage and temperature limits are tightened by the margins in Table~\ref{tab:experimental_settings}, while the predicted SOC is checked against $SOC_{\max}$. The largest candidate whose corrected predictions satisfy these bounds is selected as the executed current.}

\newrev{If no feasible current is found for any cell, the episode terminates without advancing the battery model or executing the requested action. The same model-based safety mechanism is used during training and testing.}

\subsection{Training Procedure}

\newrev{Transitions are collected in a temporary episode buffer $\mathcal{E}$ and added to the replay buffer $\mathcal{D}$ as one episode when collection ends. Each transition contains the state, availability mask, requested action, executed action, reward, next state, next availability mask, and terminal flag. Training samples batches of episodes from $\mathcal{D}$, and $|\mathcal{D}|$ denotes the number of stored episodes.}

The Q-function is evaluated using the requested action selected by the policy, while the executed action is included in the subsequent action-observation history. This allows the agent network to account for modifications made by the safety mechanism.

The overall training procedure is summarized in Algorithm~\ref{alg:QMIX}.

\begin{algorithm}[t]
\caption{MBA-RL training procedure}
\label{alg:QMIX}

\begin{algorithmic}[1]
\STATE Initialize the agent network $\theta$, mixing network $\phi$, and replay buffer $\mathcal{D}$
\STATE Set $\theta'\leftarrow\theta$ and $\phi'\leftarrow\phi$
\FOR{$episode=1$ to $N_{\mathrm{episode}}$}
    \STATE Initialize $\tau_0^i$ for all agents and observe $\mathbf{s}_0$
    \STATE \newrev{Initialize an empty episode buffer $\mathcal{E}$}
    \STATE $t\leftarrow0$, $done\leftarrow\mathrm{False}$
    \WHILE{$\neg done$}
        \STATE Compute the availability mask $\mathbf{m}_t$
        \STATE Update $\epsilon$ according to the exploration schedule
        \FOR{each agent $i$}
            \STATE Update $\tau_t^i$ using $o_t^i$ and $\tilde{a}_{t-1}^i$
            \STATE Select requested action $a_t^i$ according to \eqref{epsilon-greedy}
        \ENDFOR
        \STATE \newrev{Evaluate candidate currents using the model-based safety mechanism}
        \IF{\newrev{no feasible current exists for at least one cell}}
            \STATE \newrev{Store the terminal transition returned by the environment, with $d_t=1$, in $\mathcal{E}$}
            \STATE \newrev{End the episode without advancing the battery model}
            \STATE \textbf{break}
        \ENDIF
        \STATE \newrev{Set $\tilde{\mathbf{a}}_t$ to the feasible executed currents}
        \STATE Execute $\tilde{\mathbf{a}}_t$ and observe $r_t$, $\mathbf{s}_{t+1}$, and $done$
        \STATE Compute the next availability mask $\mathbf{m}_{t+1}$
        \STATE Store $\langle
        \mathbf{s}_t,\mathbf{m}_t,\mathbf{a}_t,
        \tilde{\mathbf{a}}_t,r_t,
        \mathbf{s}_{t+1},\mathbf{m}_{t+1},done
        \rangle$ in $\mathcal{E}$
        \STATE $t\leftarrow t+1$
    \ENDWHILE

    \IF{$\mathcal{E}\neq\varnothing$}
        \STATE \newrev{Add the collected episode $\mathcal{E}$ to $\mathcal{D}$}
    \ENDIF

    \IF{$|\mathcal{D}|\geq N_{\mathrm{batch}}$}
        \STATE Sample a batch of $N_{\mathrm{batch}}$ episodes from $\mathcal{D}$
        \STATE Compute $Q_{\mathrm{tot}}$ using \eqref{Q-value}
        \STATE Compute $y_t^{\mathrm{tot}}$ using \eqref{target Q}
        \STATE Update $\theta$ and $\phi$ using RMSprop \cite{zou2019sufficient} by minimizing
        $\left(y_t^{\mathrm{tot}}-Q_{\mathrm{tot}}\right)^2$
    \ENDIF

    \IF{$episode \bmod N_t=0$}
        \STATE $\theta'\leftarrow\theta$, $\phi'\leftarrow\phi$
    \ENDIF
\ENDFOR
\end{algorithmic}
\end{algorithm}
